\section{Learning Objectives}

After completing this chapter, readers will be able to:
\begin{itemize}
    \item Design and implement production-grade ML training pipelines
    \item Choose appropriate orchestration frameworks (Airflow, Kubeflow, MLflow)
    \item Implement comprehensive experiment tracking for reproducibility
    \item Apply hyperparameter optimization strategies at scale
    \item Build robust model validation and evaluation frameworks
    \item Design and execute A/B tests for ML model deployment
    \item Use BigQuery ML for scalable time series forecasting
    \item Implement failure handling and retry mechanisms in pipelines
    \item Apply pipeline-as-code principles for maintainability
    \item Optimize training pipelines for cost and performance
\end{itemize}

\section{Introduction: From Notebooks to Production Pipelines}

The transition from experimental notebooks to production training pipelines represents one of the most significant challenges in ML engineering. While a data scientist can train a model in a Jupyter notebook in hours, building a reliable, scalable, and maintainable training pipeline often takes weeks.

\subsection{The Training Pipeline Challenge}

Production training pipelines must address challenges that rarely appear in experimental code:

\begin{itemize}
    \item \textbf{Reproducibility:} Exact replication of experiments months or years later
    \item \textbf{Scalability:} Training on datasets that don't fit in memory
    \item \textbf{Reliability:} Handling failures, retries, and partial progress
    \item \textbf{Observability:} Tracking metrics, logs, and artifacts across runs
    \item \textbf{Cost Management:} Optimizing compute resources and storage
    \item \textbf{Version Control:} Tracking code, data, models, and configurations
    \item \textbf{Orchestration:} Coordinating multi-step workflows with dependencies
\end{itemize}

\begin{example}[The Cost of Poor Pipeline Design]
A fintech company's credit risk model degraded silently over 6 months:

\begin{itemize}
    \item Training pipeline used hardcoded dates instead of dynamic lookback
    \item Model trained on increasingly stale data as time passed
    \item No monitoring detected the data freshness issue
    \item Default rates increased 40\% before the issue was discovered
    \item Root cause: pipeline scheduled daily but used fixed date range from January
    \item Impact: \$15M in additional credit losses
\end{itemize}

The failure wasn't in the model—it was in the training pipeline design.
\end{example}

\subsection{Training Pipeline Architecture}

A production training pipeline consists of several coordinated components:

\begin{figure}[h]
\centering
\begin{tikzpicture}[
    node distance=1.8cm,
    box/.style={rectangle, draw, fill=blue!20, text width=2.3cm, text centered, minimum height=0.9cm, font=\small},
    storage/.style={cylinder, draw, fill=green!20, text width=1.8cm, text centered, minimum height=0.8cm, aspect=0.3, font=\small},
    arrow/.style={->, >=stealth, thick}
]
    % Pipeline stages
    \node[box] (data) {Data Extraction};
    \node[box, right of=data, xshift=2cm] (validate) {Validation};
    \node[box, right of=validate, xshift=2cm] (features) {Feature Engineering};
    \node[box, below of=features, yshift=-0.5cm] (train) {Model Training};
    \node[box, left of=train, xshift=-2cm] (evaluate) {Evaluation};
    \node[box, left of=evaluate, xshift=-2cm] (register) {Model Registry};

    % Storage
    \node[storage, above of=validate, yshift=0.8cm] (feature_store) {Feature Store};
    \node[storage, above of=train, yshift=0.8cm] (experiment) {Experiment Tracker};
    \node[storage, below of=register, yshift=-0.8cm] (model_store) {Model Store};

    % Orchestrator
    \node[box, below of=data, yshift=-2cm, fill=orange!30] (orchestrator) {Orchestrator};

    % Arrows - pipeline flow
    \draw[arrow] (data) -- (validate);
    \draw[arrow] (validate) -- (features);
    \draw[arrow] (features) -- (train);
    \draw[arrow] (train) -- (evaluate);
    \draw[arrow] (evaluate) -- (register);

    % Arrows - storage
    \draw[arrow, dashed] (features) -- (feature_store);
    \draw[arrow, dashed] (train) -- (experiment);
    \draw[arrow, dashed] (register) -- (model_store);

    % Orchestrator connections
    \draw[arrow, dashed] (orchestrator) -- (data);
    \draw[arrow, dashed] (orchestrator) -- (evaluate);

    % Labels
    \node[above, font=\tiny] at ($(data)!0.5!(validate)$) {Quality Gates};
    \node[above, font=\tiny] at ($(train)!0.5!(evaluate)$) {Metrics};
\end{tikzpicture}
\caption{Production Training Pipeline Architecture}
\label{fig:training_pipeline}
\end{figure}

\subsection{Pipeline Design Principles}

Effective training pipelines follow these principles:

\begin{enumerate}
    \item \textbf{Idempotency:} Running the pipeline multiple times with same inputs produces same outputs
    \item \textbf{Atomicity:} Each step either completes fully or fails cleanly
    \item \textbf{Observability:} All pipeline state, metrics, and artifacts are tracked
    \item \textbf{Configurability:} Behavior controlled by versioned configuration files
    \item \textbf{Testability:} Each component can be tested independently
    \item \textbf{Incrementality:} Only recompute what changed, reuse cached results
    \item \textbf{Fault Tolerance:} Graceful handling of failures with retry logic
\end{enumerate}

\begin{bestpractice}
\textbf{Pipeline-as-Code:} Define pipelines as version-controlled code rather than UI clicks. This enables:
\begin{itemize}
    \item Code review for pipeline changes
    \item Automated testing of pipeline logic
    \item Easy rollback to previous versions
    \item Clear audit trail of modifications
    \item Reusability across projects
\end{itemize}
\end{bestpractice}

\section{Pipeline Orchestration Frameworks}

Orchestration frameworks coordinate the execution of multi-step ML workflows, handling dependencies, scheduling, and failure recovery.

\subsection{Framework Comparison}

\begin{table}[h]
\centering
\caption{ML Orchestration Framework Comparison}
\label{tab:orchestration_comparison}
\small
\begin{tabular}{@{}lp{3cm}p{3cm}p{3cm}@{}}
\toprule
\textbf{Framework} & \textbf{Best For} & \textbf{Strengths} & \textbf{Limitations} \\
\midrule
Airflow & General workflows, batch jobs & Mature, flexible, UI & Not ML-specific, complex setup \\
Kubeflow & K8s-native ML & End-to-end ML, scalable & K8s required, steep learning curve \\
MLflow & Experiment tracking & Simple, lightweight & Limited orchestration features \\
Prefect & Modern workflows & Python-native, dynamic & Newer ecosystem \\
Metaflow & Netflix-style workflows & Data science focused & Opinionated \\
\bottomrule
\end{tabular}
\end{table}

\subsection{Apache Airflow for ML Pipelines}

Airflow remains the most widely adopted orchestration framework. Here's a production-grade ML training pipeline:

\begin{lstlisting}[style=python, caption=Production ML Training Pipeline with Airflow]
from airflow import DAG
from airflow.operators.python import PythonOperator
from airflow.providers.google.cloud.operators.bigquery import BigQueryInsertJobOperator
from airflow.providers.google.cloud.transfers.gcs_to_bigquery import GCSToBigQueryOperator
from airflow.utils.dates import days_ago
from datetime import datetime, timedelta
import mlflow
import pandas as pd
from typing import Dict, Any
import logging

logger = logging.getLogger(__name__)

# DAG configuration
default_args = {
    'owner': 'ml-team',
    'depends_on_past': False,
    'email_on_failure': True,
    'email_on_retry': False,
    'email': ['ml-alerts@company.com'],
    'retries': 2,
    'retry_delay': timedelta(minutes=5),
    'execution_timeout': timedelta(hours=2)
}

dag = DAG(
    'ml_training_pipeline',
    default_args=default_args,
    description='End-to-end ML model training pipeline',
    schedule_interval='0 2 * * *',  # Daily at 2 AM
    start_date=days_ago(1),
    catchup=False,
    max_active_runs=1,
    tags=['ml', 'training', 'production']
)

# Configuration
FEATURE_LOOKBACK_DAYS = 90
MIN_TRAINING_SAMPLES = 10000
MODEL_PERFORMANCE_THRESHOLD = 0.75
MLFLOW_TRACKING_URI = 'http://mlflow-server:5000'
EXPERIMENT_NAME = 'customer_churn_prediction'

def validate_data_quality(**context) -> Dict[str, Any]:
    """Validate data quality before training"""
    from great_expectations.data_context import DataContext

    execution_date = context['ds']
    logger.info(f"Validating data for {execution_date}")

    # Initialize Great Expectations
    ge_context = DataContext('./great_expectations')

    # Load data
    query = f"""
        SELECT *
        FROM `project.dataset.training_data`
        WHERE DATE(event_timestamp) BETWEEN
            DATE_SUB('{execution_date}', INTERVAL {FEATURE_LOOKBACK_DAYS} DAY)
            AND '{execution_date}'
    """

    from google.cloud import bigquery
    client = bigquery.Client()
    df = client.query(query).to_dataframe()

    # Check minimum sample size
    if len(df) < MIN_TRAINING_SAMPLES:
        raise ValueError(
            f"Insufficient training data: {len(df)} samples "
            f"(minimum: {MIN_TRAINING_SAMPLES})"
        )

    # Run Great Expectations validation
    batch = ge_context.get_batch({
        'datasource_name': 'pandas',
        'data_asset_name': 'training_data',
        'data': df
    }, 'training_data_expectations')

    results = batch.validate()

    if not results.success:
        failed_expectations = [
            r.expectation_config.expectation_type
            for r in results.results if not r.success
        ]
        raise ValueError(
            f"Data validation failed: {len(failed_expectations)} "
            f"expectations failed: {failed_expectations[:5]}"
        )

    # Store validation results
    validation_summary = {
        'sample_count': len(df),
        'feature_count': len(df.columns),
        'validation_success': results.success,
        'success_percent': results.statistics['success_percent']
    }

    context['task_instance'].xcom_push(
        key='validation_summary',
        value=validation_summary
    )

    logger.info(f"Validation passed: {validation_summary}")
    return validation_summary

def engineer_features(**context) -> str:
    """Feature engineering with versioning"""
    from feature_engineering import FeatureEngineer
    from feature_store import FeatureRegistry

    execution_date = context['ds']
    logger.info(f"Engineering features for {execution_date}")

    # Load raw data
    from google.cloud import bigquery
    client = bigquery.Client()

    query = f"""
        SELECT user_id, event_timestamp, event_type, value
        FROM `project.dataset.events`
        WHERE DATE(event_timestamp) BETWEEN
            DATE_SUB('{execution_date}', INTERVAL {FEATURE_LOOKBACK_DAYS} DAY)
            AND '{execution_date}'
    """

    df = client.query(query).to_dataframe()

    # Apply feature engineering
    engineer = FeatureEngineer(version='v2.1.0')
    features_df = engineer.transform(df)

    # Save features
    output_path = f'gs://ml-features/training/{execution_date}/features.parquet'
    features_df.to_parquet(output_path)

    # Update feature store
    registry = FeatureRegistry(repo_path='feature_repo/')
    registry.materialize_features(
        start_date=execution_date,
        end_date=execution_date,
        feature_views=['user_behavior_features']
    )

    logger.info(f"Engineered {len(features_df)} feature rows")

    context['task_instance'].xcom_push(
        key='features_path',
        value=output_path
    )

    return output_path

def train_model(**context) -> str:
    """Train model with MLflow tracking"""
    import mlflow
    from sklearn.ensemble import RandomForestClassifier
    from sklearn.model_selection import train_test_split
    import pandas as pd

    execution_date = context['ds']
    features_path = context['task_instance'].xcom_pull(
        task_ids='feature_engineering',
        key='features_path'
    )

    logger.info(f"Training model for {execution_date}")

    # Set MLflow tracking
    mlflow.set_tracking_uri(MLFLOW_TRACKING_URI)
    mlflow.set_experiment(EXPERIMENT_NAME)

    # Load features
    df = pd.read_parquet(features_path)

    # Prepare data
    feature_cols = [c for c in df.columns if c not in ['user_id', 'label']]
    X = df[feature_cols]
    y = df['label']

    X_train, X_test, y_train, y_test = train_test_split(
        X, y, test_size=0.2, random_state=42, stratify=y
    )

    # Start MLflow run
    with mlflow.start_run(run_name=f'training_{execution_date}') as run:

        # Log parameters
        params = {
            'n_estimators': 100,
            'max_depth': 10,
            'min_samples_split': 100,
            'class_weight': 'balanced',
            'n_jobs': -1
        }
        mlflow.log_params(params)

        # Log data info
        mlflow.log_param('training_samples', len(X_train))
        mlflow.log_param('test_samples', len(X_test))
        mlflow.log_param('feature_count', len(feature_cols))
        mlflow.log_param('execution_date', execution_date)

        # Train model
        model = RandomForestClassifier(**params)
        model.fit(X_train, y_train)

        # Evaluate
        from sklearn.metrics import (
            accuracy_score, precision_score, recall_score,
            f1_score, roc_auc_score
        )

        y_pred = model.predict(X_test)
        y_proba = model.predict_proba(X_test)[:, 1]

        metrics = {
            'accuracy': accuracy_score(y_test, y_pred),
            'precision': precision_score(y_test, y_pred),
            'recall': recall_score(y_test, y_pred),
            'f1': f1_score(y_test, y_pred),
            'roc_auc': roc_auc_score(y_test, y_proba)
        }

        mlflow.log_metrics(metrics)

        # Log feature importance
        feature_importance = pd.DataFrame({
            'feature': feature_cols,
            'importance': model.feature_importances_
        }).sort_values('importance', ascending=False)

        mlflow.log_table(
            feature_importance.head(20),
            'top_features.json'
        )

        # Log model
        mlflow.sklearn.log_model(
            model,
            'model',
            registered_model_name='churn_prediction_model'
        )

        # Check performance threshold
        if metrics['roc_auc'] < MODEL_PERFORMANCE_THRESHOLD:
            logger.warning(
                f"Model performance below threshold: "
                f"{metrics['roc_auc']:.3f} < {MODEL_PERFORMANCE_THRESHOLD}"
            )

        logger.info(f"Training complete. Metrics: {metrics}")

        # Store run info
        context['task_instance'].xcom_push(
            key='mlflow_run_id',
            value=run.info.run_id
        )
        context['task_instance'].xcom_push(
            key='model_metrics',
            value=metrics
        )

        return run.info.run_id

def evaluate_model(**context) -> Dict[str, Any]:
    """Evaluate model on holdout set"""
    import mlflow
    import pandas as pd
    from sklearn.metrics import classification_report, confusion_matrix

    run_id = context['task_instance'].xcom_pull(
        task_ids='train_model',
        key='mlflow_run_id'
    )

    logger.info(f"Evaluating model from run {run_id}")

    mlflow.set_tracking_uri(MLFLOW_TRACKING_URI)

    # Load model
    model_uri = f'runs:/{run_id}/model'
    model = mlflow.sklearn.load_model(model_uri)

    # Load holdout data (last week)
    from google.cloud import bigquery
    client = bigquery.Client()

    execution_date = context['ds']
    query = f"""
        SELECT *
        FROM `project.dataset.holdout_features`
        WHERE DATE(event_timestamp) = '{execution_date}'
    """

    df = client.query(query).to_dataframe()

    feature_cols = [c for c in df.columns if c not in ['user_id', 'label']]
    X_holdout = df[feature_cols]
    y_holdout = df['label']

    # Predictions
    y_pred = model.predict(X_holdout)
    y_proba = model.predict_proba(X_holdout)[:, 1]

    # Detailed evaluation
    from sklearn.metrics import roc_auc_score, average_precision_score

    evaluation_results = {
        'holdout_roc_auc': roc_auc_score(y_holdout, y_proba),
        'holdout_avg_precision': average_precision_score(y_holdout, y_proba),
        'holdout_samples': len(X_holdout),
        'positive_rate': y_holdout.mean()
    }

    # Log to MLflow
    with mlflow.start_run(run_id=run_id):
        mlflow.log_metrics({
            f'holdout_{k}': v for k, v in evaluation_results.items()
        })

        # Log classification report
        report = classification_report(
            y_holdout, y_pred, output_dict=True
        )
        mlflow.log_dict(report, 'classification_report.json')

        # Log confusion matrix
        cm = confusion_matrix(y_holdout, y_pred)
        mlflow.log_dict({
            'confusion_matrix': cm.tolist()
        }, 'confusion_matrix.json')

    logger.info(f"Evaluation complete: {evaluation_results}")

    # Decision: promote to staging or reject
    if evaluation_results['holdout_roc_auc'] >= MODEL_PERFORMANCE_THRESHOLD:
        context['task_instance'].xcom_push(
            key='promotion_decision',
            value='promote'
        )
    else:
        context['task_instance'].xcom_push(
            key='promotion_decision',
            value='reject'
        )
        logger.warning(
            f"Model rejected: holdout_roc_auc "
            f"{evaluation_results['holdout_roc_auc']:.3f} "
            f"< {MODEL_PERFORMANCE_THRESHOLD}"
        )

    return evaluation_results

def register_model(**context):
    """Register model to model registry if approved"""
    import mlflow
    from mlflow.tracking import MlflowClient

    promotion_decision = context['task_instance'].xcom_pull(
        task_ids='evaluate_model',
        key='promotion_decision'
    )

    if promotion_decision != 'promote':
        logger.info("Model not promoted - skipping registration")
        return

    run_id = context['task_instance'].xcom_pull(
        task_ids='train_model',
        key='mlflow_run_id'
    )

    logger.info(f"Registering model from run {run_id}")

    mlflow.set_tracking_uri(MLFLOW_TRACKING_URI)
    client = MlflowClient()

    # Get model version
    model_name = 'churn_prediction_model'
    model_uri = f'runs:/{run_id}/model'

    # Register
    model_version = mlflow.register_model(model_uri, model_name)

    # Transition to staging
    client.transition_model_version_stage(
        name=model_name,
        version=model_version.version,
        stage='Staging',
        archive_existing_versions=False
    )

    # Add description
    execution_date = context['ds']
    client.update_model_version(
        name=model_name,
        version=model_version.version,
        description=f"Trained on {execution_date}. "
                   f"Automated training pipeline run."
    )

    logger.info(
        f"Model registered: {model_name} v{model_version.version} "
        f"-> Staging"
    )

# Define tasks
validate_task = PythonOperator(
    task_id='validate_data',
    python_callable=validate_data_quality,
    dag=dag
)

feature_task = PythonOperator(
    task_id='feature_engineering',
    python_callable=engineer_features,
    dag=dag
)

train_task = PythonOperator(
    task_id='train_model',
    python_callable=train_model,
    dag=dag
)

evaluate_task = PythonOperator(
    task_id='evaluate_model',
    python_callable=evaluate_model,
    dag=dag
)

register_task = PythonOperator(
    task_id='register_model',
    python_callable=register_model,
    dag=dag
)

# Define dependencies
validate_task >> feature_task >> train_task >> evaluate_task >> register_task
\end{lstlisting}

This production pipeline demonstrates:
\begin{itemize}
    \item Data quality validation with Great Expectations
    \item Feature engineering with versioning
    \item MLflow experiment tracking
    \item Comprehensive model evaluation
    \item Conditional model promotion
    \item Error handling with retries
    \item Observable metrics at each stage
\end{itemize}

\subsection{Kubeflow Pipelines for ML}

Kubeflow Pipelines provides Kubernetes-native orchestration with first-class ML support. It excels at distributed training and GPU workloads.

\begin{lstlisting}[style=python, caption=Kubeflow Pipeline for Distributed Training]
from kfp import dsl, compiler
from kfp.v2.dsl import component, Output, Input, Dataset, Model, Metrics
from typing import NamedTuple

# Define reusable components
@component(
    base_image='python:3.9',
    packages_to_install=['pandas', 'google-cloud-bigquery', 'pyarrow']
)
def extract_data(
    project_id: str,
    dataset_id: str,
    table_id: str,
    execution_date: str,
    output_data: Output[Dataset]
) -> NamedTuple('Outputs', [('row_count', int)]):
    """Extract training data from BigQuery"""
    from google.cloud import bigquery
    import pandas as pd
    from collections import namedtuple

    client = bigquery.Client(project=project_id)

    query = f"""
        SELECT *
        FROM `{project_id}.{dataset_id}.{table_id}`
        WHERE DATE(event_timestamp) = '{execution_date}'
    """

    df = client.query(query).to_dataframe()

    # Save to component output
    df.to_parquet(output_data.path)

    outputs = namedtuple('Outputs', ['row_count'])
    return outputs(len(df))

@component(
    base_image='python:3.9',
    packages_to_install=['pandas', 'scikit-learn', 'pyarrow']
)
def preprocess_data(
    input_data: Input[Dataset],
    output_features: Output[Dataset],
    test_size: float = 0.2
) -> NamedTuple('Outputs', [('feature_count', int)]):
    """Preprocess and split data"""
    import pandas as pd
    from sklearn.model_selection import train_test_split
    from collections import namedtuple

    # Load data
    df = pd.read_parquet(input_data.path)

    # Feature engineering
    feature_cols = [c for c in df.columns if c not in ['user_id', 'label']]
    X = df[feature_cols]
    y = df['label']

    # Save processed features
    result_df = pd.DataFrame(X)
    result_df['label'] = y
    result_df.to_parquet(output_features.path)

    outputs = namedtuple('Outputs', ['feature_count'])
    return outputs(len(feature_cols))

@component(
    base_image='pytorch/pytorch:2.0.0-cuda11.7-cudnn8-runtime',
    packages_to_install=['mlflow', 'pandas', 'scikit-learn']
)
def train_model(
    input_features: Input[Dataset],
    output_model: Output[Model],
    metrics: Output[Metrics],
    n_estimators: int = 100,
    max_depth: int = 10,
    mlflow_tracking_uri: str = ''
) -> NamedTuple('Outputs', [('accuracy', float), ('roc_auc', float)]):
    """Train model with GPU support"""
    import pandas as pd
    from sklearn.ensemble import RandomForestClassifier
    from sklearn.model_selection import train_test_split
    from sklearn.metrics import accuracy_score, roc_auc_score
    import mlflow
    import pickle
    from collections import namedtuple

    # Load data
    df = pd.read_parquet(input_features.path)

    # Prepare data
    X = df.drop('label', axis=1)
    y = df['label']

    X_train, X_test, y_train, y_test = train_test_split(
        X, y, test_size=0.2, random_state=42
    )

    # Set MLflow tracking
    if mlflow_tracking_uri:
        mlflow.set_tracking_uri(mlflow_tracking_uri)

    with mlflow.start_run():
        # Train
        model = RandomForestClassifier(
            n_estimators=n_estimators,
            max_depth=max_depth,
            n_jobs=-1
        )
        model.fit(X_train, y_train)

        # Evaluate
        y_pred = model.predict(X_test)
        y_proba = model.predict_proba(X_test)[:, 1]

        acc = accuracy_score(y_test, y_pred)
        auc = roc_auc_score(y_test, y_proba)

        # Log metrics
        mlflow.log_metrics({
            'accuracy': acc,
            'roc_auc': auc
        })

        # Save model
        with open(output_model.path, 'wb') as f:
            pickle.dump(model, f)

        # Log metrics to Kubeflow
        metrics.log_metric('accuracy', acc)
        metrics.log_metric('roc_auc', auc)

    outputs = namedtuple('Outputs', ['accuracy', 'roc_auc'])
    return outputs(acc, auc)

@component(
    base_image='python:3.9',
    packages_to_install=['mlflow']
)
def deploy_model(
    input_model: Input[Model],
    model_accuracy: float,
    accuracy_threshold: float = 0.75
) -> str:
    """Deploy model if it meets threshold"""
    import mlflow
    from mlflow.tracking import MlflowClient

    if model_accuracy < accuracy_threshold:
        return f"Model rejected: accuracy {model_accuracy:.3f} < {accuracy_threshold}"

    # Transition to production
    client = MlflowClient()

    # In production, you'd actually deploy here
    # For demo, just log the decision
    return f"Model approved for deployment: accuracy {model_accuracy:.3f}"

# Define pipeline
@dsl.pipeline(
    name='ml-training-pipeline',
    description='End-to-end ML training pipeline on Kubeflow'
)
def ml_training_pipeline(
    project_id: str = 'my-project',
    dataset_id: str = 'ml_data',
    table_id: str = 'training_data',
    execution_date: str = '2024-01-15',
    n_estimators: int = 100,
    max_depth: int = 10,
    mlflow_tracking_uri: str = 'http://mlflow:5000'
):
    """Complete ML training pipeline"""

    # Extract data
    extract_task = extract_data(
        project_id=project_id,
        dataset_id=dataset_id,
        table_id=table_id,
        execution_date=execution_date
    )

    # Preprocess
    preprocess_task = preprocess_data(
        input_data=extract_task.outputs['output_data']
    )

    # Train
    train_task = train_model(
        input_features=preprocess_task.outputs['output_features'],
        n_estimators=n_estimators,
        max_depth=max_depth,
        mlflow_tracking_uri=mlflow_tracking_uri
    )

    # Deploy (conditional)
    deploy_task = deploy_model(
        input_model=train_task.outputs['output_model'],
        model_accuracy=train_task.outputs['accuracy']
    )

# Compile pipeline
if __name__ == '__main__':
    compiler.Compiler().compile(
        pipeline_func=ml_training_pipeline,
        package_path='ml_training_pipeline.yaml'
    )
\end{lstlisting}

Key advantages of Kubeflow:
\begin{itemize}
    \item \textbf{Containerization:} Each component runs in isolated container
    \item \textbf{Scalability:} Native Kubernetes scaling and GPU support
    \item \textbf{Reproducibility:} Pipeline defined as code with versioned containers
    \item \textbf{ML-Native:} Built-in support for models, datasets, and metrics
    \item \textbf{Experimentation:} Easy parameter sweeping and parallel runs
\end{itemize}

\section{Experiment Tracking and Reproducibility}

Reproducibility is fundamental to ML engineering. Every experiment must be completely reproducible months or years later.

\subsection{The Reproducibility Challenge}

Reproducing an ML experiment requires tracking:
\begin{itemize}
    \item Code version (git commit)
    \item Data version (DVC hash or timestamp)
    \item Model architecture and hyperparameters
    \item Library versions (requirements.txt)
    \item Random seeds
    \item Hardware configuration (CPU/GPU)
    \item Feature transformations
    \item Training procedure (epochs, batch size, etc.)
\end{itemize}

\subsection{MLflow for Comprehensive Tracking}

MLflow provides a complete experiment tracking solution:

\begin{lstlisting}[style=python, caption=Comprehensive Experiment Tracking with MLflow]
import mlflow
import mlflow.pytorch
from mlflow.models.signature import infer_signature
import torch
import torch.nn as nn
from torch.utils.data import DataLoader
import pandas as pd
import numpy as np
from typing import Dict, Any
import json
import git
import os

class ExperimentTracker:
    """Comprehensive experiment tracking wrapper"""

    def __init__(
        self,
        tracking_uri: str,
        experiment_name: str,
        artifact_location: str = None
    ):
        mlflow.set_tracking_uri(tracking_uri)

        # Create or get experiment
        experiment = mlflow.get_experiment_by_name(experiment_name)
        if experiment is None:
            experiment_id = mlflow.create_experiment(
                experiment_name,
                artifact_location=artifact_location
            )
        else:
            experiment_id = experiment.experiment_id

        mlflow.set_experiment(experiment_name)
        self.experiment_id = experiment_id

    def start_run(
        self,
        run_name: str,
        description: str = None,
        tags: Dict[str, str] = None
    ):
        """Start a new MLflow run with comprehensive metadata"""

        # Get git info
        repo = git.Repo(search_parent_directories=True)
        git_commit = repo.head.object.hexsha
        git_branch = repo.active_branch.name

        # Combine tags
        all_tags = {
            'git_commit': git_commit,
            'git_branch': git_branch,
            'python_version': f"{os.sys.version_info.major}.{os.sys.version_info.minor}",
            'mlflow_version': mlflow.__version__,
        }

        if tags:
            all_tags.update(tags)

        # Start run
        run = mlflow.start_run(
            run_name=run_name,
            description=description,
            tags=all_tags
        )

        return run

    def log_dataset_info(
        self,
        train_df: pd.DataFrame,
        val_df: pd.DataFrame,
        test_df: pd.DataFrame
    ):
        """Log comprehensive dataset information"""

        # Dataset sizes
        mlflow.log_param('train_samples', len(train_df))
        mlflow.log_param('val_samples', len(val_df))
        mlflow.log_param('test_samples', len(test_df))

        # Feature info
        mlflow.log_param('feature_count', len(train_df.columns))
        mlflow.log_param('features', list(train_df.columns))

        # Class distribution
        if 'label' in train_df.columns:
            class_dist = train_df['label'].value_counts().to_dict()
            mlflow.log_param('class_distribution', class_dist)

            # Class imbalance ratio
            if len(class_dist) == 2:
                imbalance = max(class_dist.values()) / min(class_dist.values())
                mlflow.log_metric('class_imbalance_ratio', imbalance)

        # Statistical summary
        summary = train_df.describe().to_dict()
        mlflow.log_dict(summary, 'dataset_summary.json')

    def log_model_architecture(
        self,
        model: nn.Module,
        input_example: torch.Tensor
    ):
        """Log model architecture and complexity"""

        # Count parameters
        total_params = sum(p.numel() for p in model.parameters())
        trainable_params = sum(
            p.numel() for p in model.parameters() if p.requires_grad
        )

        mlflow.log_param('total_parameters', total_params)
        mlflow.log_param('trainable_parameters', trainable_params)

        # Model summary
        mlflow.log_text(str(model), 'model_architecture.txt')

        # Log computational graph (requires torchviz)
        try:
            from torchviz import make_dot
            output = model(input_example)
            dot = make_dot(output, params=dict(model.named_parameters()))
            dot.render('model_graph', format='png')
            mlflow.log_artifact('model_graph.png')
        except ImportError:
            pass

    def log_training_config(
        self,
        config: Dict[str, Any]
    ):
        """Log all training hyperparameters"""

        # Flatten nested config
        def flatten_dict(d, parent_key='', sep='_'):
            items = []
            for k, v in d.items():
                new_key = f"{parent_key}{sep}{k}" if parent_key else k
                if isinstance(v, dict):
                    items.extend(flatten_dict(v, new_key, sep=sep).items())
                else:
                    items.append((new_key, v))
            return dict(items)

        flat_config = flatten_dict(config)

        # Log as parameters
        for key, value in flat_config.items():
            mlflow.log_param(key, value)

        # Also save full config as artifact
        mlflow.log_dict(config, 'full_config.json')

    def log_training_metrics(
        self,
        epoch: int,
        train_metrics: Dict[str, float],
        val_metrics: Dict[str, float]
    ):
        """Log metrics for each epoch"""

        # Log train metrics
        for name, value in train_metrics.items():
            mlflow.log_metric(f'train_{name}', value, step=epoch)

        # Log validation metrics
        for name, value in val_metrics.items():
            mlflow.log_metric(f'val_{name}', value, step=epoch)

    def log_system_metrics(self):
        """Log system and hardware metrics"""
        import psutil
        import platform

        system_info = {
            'cpu_count': psutil.cpu_count(),
            'memory_gb': psutil.virtual_memory().total / (1024**3),
            'platform': platform.platform(),
            'processor': platform.processor()
        }

        # GPU info
        if torch.cuda.is_available():
            system_info.update({
                'gpu_count': torch.cuda.device_count(),
                'gpu_name': torch.cuda.get_device_name(0),
                'cuda_version': torch.version.cuda
            })

        mlflow.log_dict(system_info, 'system_info.json')

    def log_model_artifacts(
        self,
        model: nn.Module,
        signature: mlflow.models.signature.ModelSignature = None,
        input_example: np.ndarray = None,
        requirements_file: str = 'requirements.txt'
    ):
        """Log model with all necessary artifacts"""

        # Log PyTorch model
        mlflow.pytorch.log_model(
            model,
            'model',
            signature=signature,
            input_example=input_example,
            pip_requirements=requirements_file
        )

        # Log state dict separately for flexibility
        torch.save(model.state_dict(), 'model_state.pth')
        mlflow.log_artifact('model_state.pth')

    def log_feature_importance(
        self,
        feature_names: list,
        importances: np.ndarray
    ):
        """Log feature importance analysis"""

        importance_df = pd.DataFrame({
            'feature': feature_names,
            'importance': importances
        }).sort_values('importance', ascending=False)

        # Log as table
        mlflow.log_table(importance_df, 'feature_importance.json')

        # Create and log plot
        import matplotlib.pyplot as plt

        plt.figure(figsize=(10, 8))
        plt.barh(
            importance_df['feature'].head(20),
            importance_df['importance'].head(20)
        )
        plt.xlabel('Importance')
        plt.title('Top 20 Features by Importance')
        plt.tight_layout()
        plt.savefig('feature_importance.png')
        mlflow.log_artifact('feature_importance.png')
        plt.close()

# Usage example
def train_with_tracking():
    """Example: Complete training with comprehensive tracking"""

    # Initialize tracker
    tracker = ExperimentTracker(
        tracking_uri='http://mlflow-server:5000',
        experiment_name='customer_churn_prediction'
    )

    # Start run
    with tracker.start_run(
        run_name='rf_baseline',
        description='Random Forest baseline model',
        tags={'model_type': 'random_forest', 'version': 'v1.0'}
    ):

        # Load data
        train_df = pd.read_csv('data/train.csv')
        val_df = pd.read_csv('data/val.csv')
        test_df = pd.read_csv('data/test.csv')

        # Log dataset info
        tracker.log_dataset_info(train_df, val_df, test_df)

        # Configuration
        config = {
            'model': {
                'type': 'RandomForest',
                'n_estimators': 100,
                'max_depth': 10,
                'min_samples_split': 100
            },
            'training': {
                'batch_size': 256,
                'epochs': 10,
                'optimizer': 'adam',
                'learning_rate': 0.001
            },
            'data': {
                'feature_engineering_version': 'v2.1',
                'normalization': 'standard'
            }
        }

        tracker.log_training_config(config)

        # Log system info
        tracker.log_system_metrics()

        # Train model (simplified)
        from sklearn.ensemble import RandomForestClassifier

        X_train = train_df.drop('label', axis=1)
        y_train = train_df['label']

        model = RandomForestClassifier(**config['model'])
        model.fit(X_train, y_train)

        # Log model artifacts
        tracker.log_model_artifacts(
            model,
            input_example=X_train.head(5).to_numpy()
        )

        # Log feature importance
        tracker.log_feature_importance(
            list(X_train.columns),
            model.feature_importances_
        )

        print("Training complete with comprehensive tracking")

# Compare experiments
def compare_experiments(experiment_name: str, metric: str = 'val_roc_auc'):
    """Compare all runs in an experiment"""

    mlflow.set_tracking_uri('http://mlflow-server:5000')

    # Get all runs
    runs = mlflow.search_runs(
        experiment_names=[experiment_name],
        order_by=[f"metrics.{metric} DESC"]
    )

    # Display top runs
    print(f"Top runs by {metric}:")
    print(runs[['run_id', f'metrics.{metric}', 'params.model_type']].head(10))

    return runs
\end{lstlisting}

This tracking system ensures complete reproducibility by capturing all aspects of the experiment.

\section{Hyperparameter Optimization at Scale}

Hyperparameter optimization (HPO) is critical for model performance but computationally expensive. Production systems require efficient, scalable HPO strategies.

\subsection{HPO Strategies Comparison}

\begin{table}[h]
\centering
\caption{Hyperparameter Optimization Strategies}
\label{tab:hpo_strategies}
\small
\begin{tabular}{@{}lp{3.5cm}p{2.5cm}p{3cm}@{}}
\toprule
\textbf{Method} & \textbf{Best For} & \textbf{Convergence} & \textbf{Compute Cost} \\
\midrule
Grid Search & Small parameter spaces & Guaranteed coverage & Very high \\
Random Search & Medium spaces, quick results & Better than grid & Medium \\
Bayesian Optimization & Expensive evaluations & Sample efficient & Low \\
Hyperband & Large-scale search & Fast early stopping & Medium-Low \\
Optuna & Modern, flexible & Adaptive & Medium \\
Ray Tune & Distributed, scalable & Parallel efficient & Low (amortized) \\
\bottomrule
\end{tabular}
\end{table}

\subsection{Bayesian Optimization with Optuna}

Optuna provides an efficient, Pythonic HPO framework:

\begin{lstlisting}[style=python, caption=Hyperparameter Optimization with Optuna]
import optuna
from optuna.integration.mlflow import MLflowCallback
import mlflow
from sklearn.ensemble import RandomForestClassifier
from sklearn.model_selection import cross_val_score
import pandas as pd
import numpy as np
from typing import Dict, Any
import logging

logger = logging.getLogger(__name__)

class HyperparameterOptimizer:
    """Production-grade hyperparameter optimization"""

    def __init__(
        self,
        X_train: pd.DataFrame,
        y_train: pd.Series,
        X_val: pd.DataFrame,
        y_val: pd.Series,
        mlflow_tracking_uri: str,
        experiment_name: str
    ):
        self.X_train = X_train
        self.y_train = y_train
        self.X_val = X_val
        self.y_val = y_val

        mlflow.set_tracking_uri(mlflow_tracking_uri)
        mlflow.set_experiment(experiment_name)

    def objective(self, trial: optuna.Trial) -> float:
        """Optuna objective function"""

        # Suggest hyperparameters
        params = {
            'n_estimators': trial.suggest_int('n_estimators', 50, 500),
            'max_depth': trial.suggest_int('max_depth', 3, 20),
            'min_samples_split': trial.suggest_int('min_samples_split', 2, 100),
            'min_samples_leaf': trial.suggest_int('min_samples_leaf', 1, 50),
            'max_features': trial.suggest_categorical(
                'max_features', ['sqrt', 'log2', None]
            ),
            'class_weight': trial.suggest_categorical(
                'class_weight', ['balanced', 'balanced_subsample', None]
            )
        }

        # Train model
        model = RandomForestClassifier(**params, n_jobs=-1, random_state=42)
        model.fit(self.X_train, self.y_train)

        # Evaluate on validation set
        from sklearn.metrics import roc_auc_score

        y_proba = model.predict_proba(self.X_val)[:, 1]
        score = roc_auc_score(self.y_val, y_proba)

        # Log intermediate values for pruning
        trial.report(score, step=0)

        # Handle pruning based on the intermediate value
        if trial.should_prune():
            raise optuna.TrialPruned()

        return score

    def optimize(
        self,
        n_trials: int = 100,
        timeout: int = 3600,
        n_jobs: int = -1
    ) -> optuna.Study:
        """Run hyperparameter optimization"""

        # Create study with pruning
        study = optuna.create_study(
            study_name=f'hpo_{pd.Timestamp.now().strftime("%Y%m%d_%H%M%S")}',
            direction='maximize',
            pruner=optuna.pruners.MedianPruner(
                n_startup_trials=5,
                n_warmup_steps=0,
                interval_steps=1
            ),
            sampler=optuna.samplers.TPESampler(seed=42)
        )

        # MLflow callback for tracking
        mlflow_callback = MLflowCallback(
            tracking_uri=mlflow.get_tracking_uri(),
            metric_name='val_roc_auc'
        )

        # Optimize
        study.optimize(
            self.objective,
            n_trials=n_trials,
            timeout=timeout,
            n_jobs=n_jobs,
            callbacks=[mlflow_callback],
            show_progress_bar=True
        )

        # Log best parameters
        logger.info(f"Best trial: {study.best_trial.number}")
        logger.info(f"Best value: {study.best_value:.4f}")
        logger.info(f"Best params: {study.best_params}")

        return study

    def get_best_model(self, study: optuna.Study):
        """Train final model with best parameters"""

        best_params = study.best_params

        # Train on full training set
        model = RandomForestClassifier(
            **best_params,
            n_jobs=-1,
            random_state=42
        )

        model.fit(self.X_train, self.y_train)

        return model, best_params

# Usage example
def run_hyperparameter_optimization():
    """Example: Complete HPO workflow"""

    # Load data
    train_df = pd.read_csv('data/train.csv')
    val_df = pd.read_csv('data/val.csv')

    X_train = train_df.drop('label', axis=1)
    y_train = train_df['label']
    X_val = val_df.drop('label', axis=1)
    y_val = val_df['label']

    # Initialize optimizer
    optimizer = HyperparameterOptimizer(
        X_train=X_train,
        y_train=y_train,
        X_val=X_val,
        y_val=y_val,
        mlflow_tracking_uri='http://mlflow-server:5000',
        experiment_name='churn_prediction_hpo'
    )

    # Run optimization
    study = optimizer.optimize(
        n_trials=100,
        timeout=3600,
        n_jobs=4
    )

    # Get best model
    best_model, best_params = optimizer.get_best_model(study)

    # Log final model
    with mlflow.start_run(run_name='best_model_from_hpo'):
        mlflow.log_params(best_params)
        mlflow.log_metric('val_roc_auc', study.best_value)
        mlflow.sklearn.log_model(best_model, 'model')

    # Visualization
    optuna.visualization.plot_optimization_history(study).show()
    optuna.visualization.plot_param_importances(study).show()

    return best_model, study
\end{lstlisting}

\subsection{Distributed HPO with Ray Tune}

For large-scale HPO, Ray Tune provides distributed, fault-tolerant optimization:

\begin{lstlisting}[style=python, caption=Distributed Hyperparameter Optimization with Ray Tune]
from ray import tune
from ray.tune import CLIReporter
from ray.tune.schedulers import ASHAScheduler
from ray.tune.integration.mlflow import mlflow_mixin
import mlflow
from typing import Dict

@mlflow_mixin
def train_model_tune(config: Dict):
    """Training function for Ray Tune"""
    import pandas as pd
    from sklearn.ensemble import RandomForestClassifier
    from sklearn.metrics import roc_auc_score

    # Load data
    train_df = pd.read_csv('data/train.csv')
    val_df = pd.read_csv('data/val.csv')

    X_train = train_df.drop('label', axis=1)
    y_train = train_df['label']
    X_val = val_df.drop('label', axis=1)
    y_val = val_df['label']

    # Train model
    model = RandomForestClassifier(
        n_estimators=config['n_estimators'],
        max_depth=config['max_depth'],
        min_samples_split=config['min_samples_split'],
        min_samples_leaf=config['min_samples_leaf'],
        n_jobs=1,  # Ray handles parallelism
        random_state=42
    )

    model.fit(X_train, y_train)

    # Evaluate
    y_proba = model.predict_proba(X_val)[:, 1]
    score = roc_auc_score(y_val, y_proba)

    # Report to Ray Tune (and MLflow via mixin)
    tune.report(val_roc_auc=score)

def run_distributed_hpo():
    """Run distributed HPO with Ray Tune"""

    # Define search space
    config = {
        'n_estimators': tune.randint(50, 500),
        'max_depth': tune.randint(3, 20),
        'min_samples_split': tune.randint(2, 100),
        'min_samples_leaf': tune.randint(1, 50)
    }

    # Configure ASHA scheduler for early stopping
    scheduler = ASHAScheduler(
        max_t=100,
        grace_period=10,
        reduction_factor=2
    )

    # Reporter for progress
    reporter = CLIReporter(
        parameter_columns=['n_estimators', 'max_depth'],
        metric_columns=['val_roc_auc']
    )

    # Run tuning
    analysis = tune.run(
        train_model_tune,
        config=config,
        metric='val_roc_auc',
        mode='max',
        num_samples=100,
        scheduler=scheduler,
        progress_reporter=reporter,
        resources_per_trial={'cpu': 4},
        mlflow={'experiment_name': 'ray_tune_hpo'}
    )

    # Get best config
    best_config = analysis.get_best_config(metric='val_roc_auc', mode='max')
    print(f"Best config: {best_config}")

    return analysis
\end{lstlisting}

\section{Model Validation and Evaluation Frameworks}

Robust model validation goes beyond simple train/test splits. Production systems require comprehensive evaluation frameworks.

\subsection{Comprehensive Validation Strategy}

\begin{lstlisting}[style=python, caption=Production Model Validation Framework]
from sklearn.model_selection import (
    StratifiedKFold, TimeSeriesSplit, cross_validate
)
from sklearn.metrics import (
    accuracy_score, precision_score, recall_score, f1_score,
    roc_auc_score, average_precision_score, confusion_matrix,
    classification_report
)
import pandas as pd
import numpy as np
from typing import Dict, List, Any, Callable
import matplotlib.pyplot as plt
import seaborn as sns
import logging

logger = logging.getLogger(__name__)

class ModelValidator:
    """Comprehensive model validation framework"""

    def __init__(self, model, X, y, is_time_series: bool = False):
        self.model = model
        self.X = X
        self.y = y
        self.is_time_series = is_time_series

    def cross_validate_comprehensive(
        self,
        n_splits: int = 5
    ) -> Dict[str, Any]:
        """Comprehensive cross-validation with multiple metrics"""

        # Choose appropriate CV strategy
        if self.is_time_series:
            cv = TimeSeriesSplit(n_splits=n_splits)
            logger.info(f"Using TimeSeriesSplit with {n_splits} splits")
        else:
            cv = StratifiedKFold(n_splits=n_splits, shuffle=True, random_state=42)
            logger.info(f"Using StratifiedKFold with {n_splits} splits")

        # Define scoring metrics
        scoring = {
            'accuracy': 'accuracy',
            'precision': 'precision',
            'recall': 'recall',
            'f1': 'f1',
            'roc_auc': 'roc_auc',
            'average_precision': 'average_precision'
        }

        # Run cross-validation
        cv_results = cross_validate(
            self.model,
            self.X,
            self.y,
            cv=cv,
            scoring=scoring,
            return_train_score=True,
            return_estimator=True,
            n_jobs=-1
        )

        # Aggregate results
        results_summary = {}

        for metric in scoring.keys():
            test_scores = cv_results[f'test_{metric}']
            train_scores = cv_results[f'train_{metric}']

            results_summary[metric] = {
                'test_mean': np.mean(test_scores),
                'test_std': np.std(test_scores),
                'test_scores': test_scores.tolist(),
                'train_mean': np.mean(train_scores),
                'train_std': np.std(train_scores),
                'overfitting_gap': np.mean(train_scores) - np.mean(test_scores)
            }

        # Check for overfitting
        for metric, scores in results_summary.items():
            if scores['overfitting_gap'] > 0.1:
                logger.warning(
                    f"Potential overfitting detected for {metric}: "
                    f"gap = {scores['overfitting_gap']:.3f}"
                )

        results_summary['cv_estimators'] = cv_results['estimator']

        return results_summary

    def evaluate_thresholds(
        self,
        X_test: pd.DataFrame,
        y_test: pd.Series,
        thresholds: List[float] = None
    ) -> pd.DataFrame:
        """Evaluate model at different decision thresholds"""

        if thresholds is None:
            thresholds = np.arange(0.1, 1.0, 0.05)

        # Get predicted probabilities
        y_proba = self.model.predict_proba(X_test)[:, 1]

        results = []

        for threshold in thresholds:
            y_pred = (y_proba >= threshold).astype(int)

            # Calculate metrics
            precision = precision_score(y_test, y_pred, zero_division=0)
            recall = recall_score(y_test, y_pred, zero_division=0)
            f1 = f1_score(y_test, y_pred, zero_division=0)

            results.append({
                'threshold': threshold,
                'precision': precision,
                'recall': recall,
                'f1': f1
            })

        results_df = pd.DataFrame(results)

        # Find optimal threshold (maximize F1)
        optimal_idx = results_df['f1'].idxmax()
        optimal_threshold = results_df.loc[optimal_idx, 'threshold']

        logger.info(
            f"Optimal threshold: {optimal_threshold:.2f} "
            f"(F1: {results_df.loc[optimal_idx, 'f1']:.3f})"
        )

        return results_df

    def evaluate_fairness(
        self,
        X_test: pd.DataFrame,
        y_test: pd.Series,
        sensitive_feature: str
    ) -> Dict[str, Any]:
        """Evaluate model fairness across groups"""

        y_pred = self.model.predict(X_test)

        # Get unique groups
        groups = X_test[sensitive_feature].unique()

        fairness_metrics = {}

        for group in groups:
            mask = X_test[sensitive_feature] == group

            group_metrics = {
                'size': mask.sum(),
                'accuracy': accuracy_score(y_test[mask], y_pred[mask]),
                'precision': precision_score(y_test[mask], y_pred[mask]),
                'recall': recall_score(y_test[mask], y_pred[mask]),
                'f1': f1_score(y_test[mask], y_pred[mask])
            }

            fairness_metrics[str(group)] = group_metrics

        # Calculate disparity
        accuracies = [m['accuracy'] for m in fairness_metrics.values()]
        disparity = max(accuracies) - min(accuracies)

        fairness_metrics['accuracy_disparity'] = disparity

        if disparity > 0.05:
            logger.warning(
                f"Fairness concern: accuracy disparity = {disparity:.3f}"
            )

        return fairness_metrics

    def generate_report(
        self,
        X_test: pd.DataFrame,
        y_test: pd.Series,
        output_dir: str = 'validation_reports'
    ):
        """Generate comprehensive validation report"""
        import os
        os.makedirs(output_dir, exist_ok=True)

        # Predictions
        y_pred = self.model.predict(X_test)
        y_proba = self.model.predict_proba(X_test)[:, 1]

        # Classification report
        report = classification_report(y_test, y_pred, output_dict=True)
        report_df = pd.DataFrame(report).transpose()
        report_df.to_csv(f'{output_dir}/classification_report.csv')

        # Confusion matrix
        cm = confusion_matrix(y_test, y_pred)

        plt.figure(figsize=(8, 6))
        sns.heatmap(cm, annot=True, fmt='d', cmap='Blues')
        plt.title('Confusion Matrix')
        plt.ylabel('True Label')
        plt.xlabel('Predicted Label')
        plt.savefig(f'{output_dir}/confusion_matrix.png')
        plt.close()

        # ROC curve
        from sklearn.metrics import roc_curve, auc

        fpr, tpr, thresholds = roc_curve(y_test, y_proba)
        roc_auc = auc(fpr, tpr)

        plt.figure(figsize=(8, 6))
        plt.plot(fpr, tpr, label=f'ROC curve (AUC = {roc_auc:.3f})')
        plt.plot([0, 1], [0, 1], 'k--', label='Random')
        plt.xlabel('False Positive Rate')
        plt.ylabel('True Positive Rate')
        plt.title('ROC Curve')
        plt.legend()
        plt.savefig(f'{output_dir}/roc_curve.png')
        plt.close()

        # Precision-Recall curve
        from sklearn.metrics import precision_recall_curve

        precision, recall, _ = precision_recall_curve(y_test, y_proba)
        avg_precision = average_precision_score(y_test, y_proba)

        plt.figure(figsize=(8, 6))
        plt.plot(recall, precision, label=f'AP = {avg_precision:.3f}')
        plt.xlabel('Recall')
        plt.ylabel('Precision')
        plt.title('Precision-Recall Curve')
        plt.legend()
        plt.savefig(f'{output_dir}/precision_recall_curve.png')
        plt.close()

        logger.info(f"Validation report saved to {output_dir}/")

# Usage
def validate_model_comprehensive():
    """Example: Comprehensive model validation"""

    from sklearn.ensemble import RandomForestClassifier

    # Load data
    train_df = pd.read_csv('data/train.csv')
    test_df = pd.read_csv('data/test.csv')

    X_train = train_df.drop('label', axis=1)
    y_train = train_df['label']
    X_test = test_df.drop('label', axis=1)
    y_test = test_df['label']

    # Train model
    model = RandomForestClassifier(n_estimators=100, random_state=42)
    model.fit(X_train, y_train)

    # Initialize validator
    validator = ModelValidator(model, X_train, y_train)

    # Cross-validation
    cv_results = validator.cross_validate_comprehensive(n_splits=5)

    # Threshold analysis
    threshold_results = validator.evaluate_thresholds(X_test, y_test)

    # Fairness evaluation
    fairness_results = validator.evaluate_fairness(
        X_test, y_test, sensitive_feature='age_group'
    )

    # Generate report
    validator.generate_report(X_test, y_test)

    return cv_results, threshold_results, fairness_results
\end{lstlisting}

\section{A/B Testing for ML Models}

A/B testing ML models in production requires careful experimental design and statistical rigor.

\subsection{Model A/B Test Design}

\begin{lstlisting}[style=python, caption=A/B Testing Framework for ML Models]
import hashlib
from typing import Dict, Any, Optional
from dataclasses import dataclass
from datetime import datetime
import pandas as pd
import numpy as np
from scipy import stats
import logging

logger = logging.getLogger(__name__)

@dataclass
class ABTestConfig:
    """Configuration for model A/B test"""
    experiment_name: str
    control_model_version: str
    treatment_model_version: str
    traffic_allocation: float = 0.5  # 50/50 split
    minimum_sample_size: int = 1000
    significance_level: float = 0.05
    minimum_detectable_effect: float = 0.02  # 2% relative change
    start_date: datetime = None
    end_date: datetime = None

class ModelABTest:
    """Production A/B testing framework for ML models"""

    def __init__(self, config: ABTestConfig):
        self.config = config
        self.results = {'control': [], 'treatment': []}

    def assign_variant(self, user_id: str) -> str:
        """Consistent user assignment using hash-based randomization"""

        # Create deterministic hash
        hash_value = int(
            hashlib.md5(
                f"{user_id}_{self.config.experiment_name}".encode()
            ).hexdigest(),
            16
        )

        # Convert to probability [0, 1)
        prob = (hash_value % 10000) / 10000

        if prob < self.config.traffic_allocation:
            return 'treatment'
        else:
            return 'control'

    def get_model_version(self, user_id: str) -> str:
        """Get model version for user"""

        variant = self.assign_variant(user_id)

        if variant == 'treatment':
            return self.config.treatment_model_version
        else:
            return self.config.control_model_version

    def log_prediction(
        self,
        user_id: str,
        prediction: float,
        actual: Optional[float] = None
    ):
        """Log prediction result for A/B test"""

        variant = self.assign_variant(user_id)

        result = {
            'user_id': user_id,
            'timestamp': datetime.now(),
            'prediction': prediction,
            'actual': actual,
            'variant': variant
        }

        self.results[variant].append(result)

    def calculate_sample_size(
        self,
        baseline_conversion: float,
        minimum_detectable_effect: float = None
    ) -> int:
        """Calculate required sample size for statistical power"""

        if minimum_detectable_effect is None:
            minimum_detectable_effect = self.config.minimum_detectable_effect

        # Use statsmodels for power analysis
        from statsmodels.stats.power import zt_ind_solve_power

        # Calculate required sample size per variant
        n_per_variant = zt_ind_solve_power(
            effect_size=minimum_detectable_effect,
            alpha=self.config.significance_level,
            power=0.8,  # 80% power
            ratio=1.0,  # Equal sample sizes
            alternative='two-sided'
        )

        total_sample_size = int(np.ceil(n_per_variant * 2))

        logger.info(
            f"Required sample size: {total_sample_size} "
            f"({n_per_variant:.0f} per variant)"
        )

        return total_sample_size

    def analyze_results(
        self,
        metric: str = 'accuracy'
    ) -> Dict[str, Any]:
        """Analyze A/B test results"""

        # Convert results to DataFrames
        control_df = pd.DataFrame(self.results['control'])
        treatment_df = pd.DataFrame(self.results['treatment'])

        # Check minimum sample size
        if len(control_df) < self.config.minimum_sample_size:
            logger.warning(
                f"Insufficient control samples: {len(control_df)} "
                f"< {self.config.minimum_sample_size}"
            )
            return {'status': 'insufficient_data'}

        if len(treatment_df) < self.config.minimum_sample_size:
            logger.warning(
                f"Insufficient treatment samples: {len(treatment_df)} "
                f"< {self.config.minimum_sample_size}"
            )
            return {'status': 'insufficient_data'}

        # Calculate metric
        if metric == 'accuracy':
            control_metric = (
                control_df['prediction'] == control_df['actual']
            ).mean()
            treatment_metric = (
                treatment_df['prediction'] == treatment_df['actual']
            ).mean()

        elif metric == 'auc':
            from sklearn.metrics import roc_auc_score
            control_metric = roc_auc_score(
                control_df['actual'],
                control_df['prediction']
            )
            treatment_metric = roc_auc_score(
                treatment_df['actual'],
                treatment_df['prediction']
            )

        # Calculate absolute and relative difference
        absolute_diff = treatment_metric - control_metric
        relative_diff = absolute_diff / control_metric

        # Statistical test (two-sample t-test)
        t_stat, p_value = stats.ttest_ind(
            treatment_df['prediction'],
            control_df['prediction'],
            equal_var=False
        )

        # Confidence interval
        se = np.sqrt(
            control_metric * (1 - control_metric) / len(control_df) +
            treatment_metric * (1 - treatment_metric) / len(treatment_df)
        )

        ci_lower = absolute_diff - 1.96 * se
        ci_upper = absolute_diff + 1.96 * se

        # Decision
        is_significant = p_value < self.config.significance_level
        recommend_treatment = is_significant and absolute_diff > 0

        results = {
            'status': 'complete',
            'control': {
                'samples': len(control_df),
                'metric': control_metric
            },
            'treatment': {
                'samples': len(treatment_df),
                'metric': treatment_metric
            },
            'difference': {
                'absolute': absolute_diff,
                'relative': relative_diff,
                'ci_lower': ci_lower,
                'ci_upper': ci_upper
            },
            'statistics': {
                't_statistic': t_stat,
                'p_value': p_value,
                'is_significant': is_significant
            },
            'recommendation': {
                'deploy_treatment': recommend_treatment,
                'reason': self._get_recommendation_reason(
                    is_significant, absolute_diff, p_value
                )
            }
        }

        logger.info(
            f"A/B Test Results: "
            f"Control={control_metric:.4f}, "
            f"Treatment={treatment_metric:.4f}, "
            f"Diff={relative_diff:.2%}, "
            f"p-value={p_value:.4f}"
        )

        return results

    def _get_recommendation_reason(
        self,
        is_significant: bool,
        diff: float,
        p_value: float
    ) -> str:
        """Generate recommendation explanation"""

        if not is_significant:
            return f"No significant difference (p={p_value:.4f})"

        if diff > 0:
            return f"Treatment significantly better (p={p_value:.4f}, improvement={diff:.2%})"
        else:
            return f"Treatment significantly worse (p={p_value:.4f}, degradation={diff:.2%})"

# Usage example
def run_model_ab_test():
    """Example: Production A/B test for model deployment"""

    # Configure test
    config = ABTestConfig(
        experiment_name='churn_model_v2_vs_v1',
        control_model_version='v1.0',
        treatment_model_version='v2.0',
        traffic_allocation=0.5,
        minimum_sample_size=1000,
        significance_level=0.05,
        minimum_detectable_effect=0.02
    )

    # Initialize test
    ab_test = ModelABTest(config)

    # Calculate required sample size
    required_n = ab_test.calculate_sample_size(
        baseline_conversion=0.75  # Current accuracy
    )

    # In production: log predictions as they happen
    # Here: simulated data
    for user_id in range(5000):
        # Get assigned model version
        model_version = ab_test.get_model_version(f'user_{user_id}')

        # Make prediction (simulated)
        prediction = np.random.choice([0, 1], p=[0.3, 0.7])
        actual = np.random.choice([0, 1], p=[0.3, 0.7])

        # Log result
        ab_test.log_prediction(
            user_id=f'user_{user_id}',
            prediction=prediction,
            actual=actual
        )

    # Analyze results
    results = ab_test.analyze_results(metric='accuracy')

    # Make deployment decision
    if results.get('recommendation', {}).get('deploy_treatment'):
        logger.info("Recommendation: Deploy treatment model to 100% traffic")
    else:
        logger.info("Recommendation: Keep control model")

    return results
\end{lstlisting}

\section{Case Study: BigQuery ML ARIMA\_PLUS for Time Series Forecasting}

BigQuery ML enables scalable time series forecasting without infrastructure management. This case study demonstrates production time series pipelines.

\subsection{Business Context}

An e-commerce company needs to forecast daily sales for inventory planning. Requirements:
\begin{itemize}
    \item 7-day ahead forecast updated daily
    \item Incorporate seasonality and holidays
    \item Handle multiple product categories
    \item Minimal infrastructure overhead
\end{itemize}

\subsection{Implementation with BigQuery ML}

\begin{lstlisting}[style=sql, caption=BigQuery ML ARIMA\_PLUS Time Series Forecast]
-- Step 1: Prepare training data
CREATE OR REPLACE TABLE `project.dataset.sales_training_data` AS
SELECT
    DATE(order_timestamp) AS date,
    product_category,
    SUM(sales_amount) AS daily_sales,
    COUNT(DISTINCT order_id) AS order_count
FROM `project.dataset.orders`
WHERE DATE(order_timestamp) BETWEEN '2022-01-01' AND '2024-01-31'
GROUP BY date, product_category
ORDER BY product_category, date;

-- Step 2: Create ARIMA_PLUS model
CREATE OR REPLACE MODEL `project.dataset.sales_forecast_model`
OPTIONS(
    model_type='ARIMA_PLUS',
    time_series_timestamp_col='date',
    time_series_data_col='daily_sales',
    time_series_id_col='product_category',
    holiday_region='US',  -- Include US holidays
    auto_arima=TRUE,
    data_frequency='DAILY',
    decompose_time_series=TRUE
) AS
SELECT
    date,
    product_category,
    daily_sales
FROM `project.dataset.sales_training_data`;

-- Step 3: Evaluate model
SELECT
    *
FROM ML.ARIMA_EVALUATE(
    MODEL `project.dataset.sales_forecast_model`,
    STRUCT(TRUE AS show_all_candidate_models)
);

-- Step 4: Generate 7-day forecast
CREATE OR REPLACE TABLE `project.dataset.sales_forecast` AS
SELECT
    *
FROM ML.FORECAST(
    MODEL `project.dataset.sales_forecast_model`,
    STRUCT(
        7 AS horizon,  -- 7 days ahead
        0.95 AS confidence_level  -- 95% confidence intervals
    )
);

-- Step 5: Explain forecast components
SELECT
    product_category,
    date,
    daily_sales_forecast,
    standard_error,
    confidence_level,
    prediction_interval_lower_bound,
    prediction_interval_upper_bound
FROM `project.dataset.sales_forecast`
ORDER BY product_category, date;

-- Step 6: Detect anomalies
CREATE OR REPLACE TABLE `project.dataset.sales_anomalies` AS
SELECT
    *
FROM ML.DETECT_ANOMALIES(
    MODEL `project.dataset.sales_forecast_model`,
    STRUCT(0.95 AS anomaly_prob_threshold)
);
\end{lstlisting}

\subsection{Production Pipeline with Airflow}

\begin{lstlisting}[style=python, caption=Production BigQuery ML Pipeline]
from airflow import DAG
from airflow.providers.google.cloud.operators.bigquery import BigQueryInsertJobOperator
from airflow.operators.python import PythonOperator
from datetime import datetime, timedelta
import pandas as pd
from google.cloud import bigquery

default_args = {
    'owner': 'ml-team',
    'depends_on_past': False,
    'email_on_failure': True,
    'retries': 2
}

dag = DAG(
    'bigquery_ml_arima_forecast',
    default_args=default_args,
    description='Daily sales forecasting with BigQuery ML',
    schedule_interval='0 3 * * *',  # Daily at 3 AM
    start_date=datetime(2024, 1, 1),
    catchup=False
)

# Task 1: Retrain model (weekly)
retrain_model = BigQueryInsertJobOperator(
    task_id='retrain_arima_model',
    configuration={
        'query': {
            'query': '''
                CREATE OR REPLACE MODEL `project.dataset.sales_forecast_model`
                OPTIONS(
                    model_type='ARIMA_PLUS',
                    time_series_timestamp_col='date',
                    time_series_data_col='daily_sales',
                    time_series_id_col='product_category',
                    holiday_region='US',
                    auto_arima=TRUE
                ) AS
                SELECT date, product_category, daily_sales
                FROM `project.dataset.sales_training_data`
                WHERE date >= DATE_SUB(CURRENT_DATE(), INTERVAL 730 DAY)
            ''',
            'useLegacySql': False
        }
    },
    dag=dag
)

# Task 2: Generate forecast
generate_forecast = BigQueryInsertJobOperator(
    task_id='generate_forecast',
    configuration={
        'query': {
            'query': '''
                CREATE OR REPLACE TABLE `project.dataset.sales_forecast` AS
                SELECT *
                FROM ML.FORECAST(
                    MODEL `project.dataset.sales_forecast_model`,
                    STRUCT(7 AS horizon, 0.95 AS confidence_level)
                )
            ''',
            'useLegacySql': False
        }
    },
    dag=dag
)

# Task 3: Validate forecast
def validate_forecast(**context):
    """Validate forecast reasonableness"""
    client = bigquery.Client()

    # Get forecast
    query = '''
        SELECT
            product_category,
            AVG(daily_sales_forecast) as avg_forecast,
            MAX(daily_sales_forecast) as max_forecast,
            MIN(daily_sales_forecast) as min_forecast
        FROM `project.dataset.sales_forecast`
        GROUP BY product_category
    '''

    forecast_df = client.query(query).to_dataframe()

    # Get historical average
    hist_query = '''
        SELECT
            product_category,
            AVG(daily_sales) as avg_historical
        FROM `project.dataset.sales_training_data`
        WHERE date >= DATE_SUB(CURRENT_DATE(), INTERVAL 90 DAY)
        GROUP BY product_category
    '''

    hist_df = client.query(hist_query).to_dataframe()

    # Merge and validate
    merged = forecast_df.merge(hist_df, on='product_category')

    # Check if forecast is within 3x historical average
    merged['is_reasonable'] = (
        (merged['avg_forecast'] <= merged['avg_historical'] * 3) &
        (merged['avg_forecast'] >= merged['avg_historical'] * 0.3)
    )

    if not merged['is_reasonable'].all():
        unreasonable = merged[~merged['is_reasonable']]
        raise ValueError(
            f"Unreasonable forecast detected for categories: "
            f"{unreasonable['product_category'].tolist()}"
        )

    print("Forecast validation passed")

validate_task = PythonOperator(
    task_id='validate_forecast',
    python_callable=validate_forecast,
    dag=dag
)

# Task 4: Export to data warehouse
export_forecast = BigQueryInsertJobOperator(
    task_id='export_forecast',
    configuration={
        'query': {
            'query': '''
                INSERT INTO `project.warehouse.daily_forecasts`
                SELECT
                    CURRENT_TIMESTAMP() as generated_at,
                    product_category,
                    date as forecast_date,
                    daily_sales_forecast,
                    prediction_interval_lower_bound,
                    prediction_interval_upper_bound
                FROM `project.dataset.sales_forecast`
            ''',
            'useLegacySql': False
        }
    },
    dag=dag
)

# Dependencies
retrain_model >> generate_forecast >> validate_task >> export_forecast
\end{lstlisting}

Key advantages of BigQuery ML for time series:
\begin{itemize}
    \item \textbf{Serverless:} No infrastructure management
    \item \textbf{Scalable:} Handles millions of time series
    \item \textbf{SQL-Native:} Familiar interface for data teams
    \item \textbf{Auto-tuning:} Automatic ARIMA parameter selection
    \item \textbf{Seasonality:} Built-in holiday and seasonal effects
    \item \textbf{Anomaly Detection:} Integrated anomaly detection
\end{itemize}

\section{Failure Handling and Retry Mechanisms}

Production pipelines must handle failures gracefully and recover automatically.

\subsection{Failure Patterns and Solutions}

\begin{table}[h]
\centering
\caption{Common ML Pipeline Failures}
\label{tab:pipeline_failures}
\small
\begin{tabular}{@{}lp{4cm}p{4.5cm}@{}}
\toprule
\textbf{Failure Type} & \textbf{Cause} & \textbf{Solution} \\
\midrule
Data unavailability & Upstream pipeline delayed & Retry with exponential backoff \\
OOM errors & Insufficient memory & Batch processing, distributed compute \\
Network timeouts & Temporary connectivity & Automatic retry with circuit breaker \\
Invalid data & Schema changes & Data validation gates \\
Resource exhaustion & High concurrent load & Rate limiting, queue management \\
Dependency failure & External service down & Fallback to cached results \\
\bottomrule
\end{tabular}
\end{table}

\subsection{Robust Failure Handling}

\begin{lstlisting}[style=python, caption=Production-Grade Failure Handling]
from functools import wraps
import time
import logging
from typing import Callable, Any, Optional
from dataclasses import dataclass
import traceback

logger = logging.getLogger(__name__)

@dataclass
class RetryConfig:
    """Retry configuration"""
    max_retries: int = 3
    base_delay: float = 1.0  # seconds
    max_delay: float = 60.0
    exponential_base: float = 2.0
    jitter: bool = True

def retry_with_backoff(config: RetryConfig = None):
    """Decorator for retry with exponential backoff"""

    if config is None:
        config = RetryConfig()

    def decorator(func: Callable) -> Callable:
        @wraps(func)
        def wrapper(*args, **kwargs) -> Any:
            last_exception = None

            for attempt in range(config.max_retries):
                try:
                    return func(*args, **kwargs)

                except Exception as e:
                    last_exception = e

                    if attempt == config.max_retries - 1:
                        logger.error(
                            f"Function {func.__name__} failed after "
                            f"{config.max_retries} attempts: {e}"
                        )
                        raise

                    # Calculate backoff delay
                    delay = min(
                        config.base_delay * (config.exponential_base ** attempt),
                        config.max_delay
                    )

                    # Add jitter
                    if config.jitter:
                        import random
                        delay *= (0.5 + random.random())

                    logger.warning(
                        f"Attempt {attempt + 1}/{config.max_retries} failed "
                        f"for {func.__name__}: {e}. "
                        f"Retrying in {delay:.1f}s..."
                    )

                    time.sleep(delay)

            raise last_exception

        return wrapper
    return decorator

class CircuitBreaker:
    """Circuit breaker pattern for external dependencies"""

    def __init__(
        self,
        failure_threshold: int = 5,
        timeout: float = 60.0,
        expected_exception: type = Exception
    ):
        self.failure_threshold = failure_threshold
        self.timeout = timeout
        self.expected_exception = expected_exception

        self.failure_count = 0
        self.last_failure_time = None
        self.state = 'closed'  # closed, open, half_open

    def call(self, func: Callable, *args, **kwargs) -> Any:
        """Execute function with circuit breaker protection"""

        if self.state == 'open':
            if time.time() - self.last_failure_time > self.timeout:
                self.state = 'half_open'
                logger.info("Circuit breaker entering half-open state")
            else:
                raise Exception("Circuit breaker is OPEN")

        try:
            result = func(*args, **kwargs)

            # Success - reset circuit breaker
            if self.state == 'half_open':
                self.state = 'closed'
                self.failure_count = 0
                logger.info("Circuit breaker closed")

            return result

        except self.expected_exception as e:
            self.failure_count += 1
            self.last_failure_time = time.time()

            if self.failure_count >= self.failure_threshold:
                self.state = 'open'
                logger.error(
                    f"Circuit breaker opened after {self.failure_count} failures"
                )

            raise

class PipelineCheckpoint:
    """Checkpoint mechanism for long-running pipelines"""

    def __init__(self, checkpoint_dir: str):
        self.checkpoint_dir = checkpoint_dir
        import os
        os.makedirs(checkpoint_dir, exist_ok=True)

    def save(self, step: str, data: Any):
        """Save checkpoint"""
        import pickle

        checkpoint_path = f"{self.checkpoint_dir}/{step}.pkl"
        with open(checkpoint_path, 'wb') as f:
            pickle.dump(data, f)

        logger.info(f"Checkpoint saved: {step}")

    def load(self, step: str) -> Optional[Any]:
        """Load checkpoint if exists"""
        import pickle
        import os

        checkpoint_path = f"{self.checkpoint_dir}/{step}.pkl"

        if os.path.exists(checkpoint_path):
            with open(checkpoint_path, 'rb') as f:
                data = pickle.load(f)

            logger.info(f"Checkpoint loaded: {step}")
            return data

        return None

    def exists(self, step: str) -> bool:
        """Check if checkpoint exists"""
        import os
        checkpoint_path = f"{self.checkpoint_dir}/{step}.pkl"
        return os.path.exists(checkpoint_path)

# Usage example
@retry_with_backoff(RetryConfig(max_retries=3, base_delay=2.0))
def train_model_with_retry(data_path: str):
    """Train model with automatic retry"""
    import pandas as pd

    # This might fail due to network issues
    df = pd.read_csv(data_path)

    # Train model
    from sklearn.ensemble import RandomForestClassifier

    X = df.drop('label', axis=1)
    y = df['label']

    model = RandomForestClassifier()
    model.fit(X, y)

    return model

def resilient_pipeline():
    """Example: Pipeline with checkpointing and retry"""

    checkpoint = PipelineCheckpoint('./checkpoints')

    # Step 1: Data loading (with checkpoint)
    if checkpoint.exists('data'):
        data = checkpoint.load('data')
    else:
        data = train_model_with_retry('data/train.csv')
        checkpoint.save('data', data)

    # Step 2: Training (with circuit breaker for external service)
    circuit_breaker = CircuitBreaker(failure_threshold=3, timeout=60)

    try:
        model = circuit_breaker.call(train_model_with_retry, 'data/train.csv')
        checkpoint.save('model', model)
    except Exception as e:
        logger.error(f"Training failed: {e}")
        raise

    return model
\end{lstlisting}

\section{Chapter Summary}

This chapter covered production training pipeline design:

\begin{itemize}
    \item \textbf{Orchestration:} Airflow for general workflows, Kubeflow for K8s-native ML
    \item \textbf{Experiment Tracking:} MLflow for comprehensive reproducibility
    \item \textbf{Hyperparameter Optimization:} Optuna and Ray Tune for efficient search
    \item \textbf{Model Validation:} Cross-validation, threshold analysis, fairness evaluation
    \item \textbf{A/B Testing:} Statistical frameworks for model deployment decisions
    \item \textbf{BigQuery ML:} Serverless time series forecasting at scale
    \item \textbf{Failure Handling:} Retry mechanisms, circuit breakers, checkpointing
\end{itemize}

Key takeaways:

\begin{enumerate}
    \item Define pipelines as code for version control and reproducibility
    \item Track every aspect of experiments—code, data, config, environment
    \item Use Bayesian optimization for expensive model training
    \item Validate models comprehensively beyond simple accuracy metrics
    \item Design A/B tests with proper statistical power
    \item Choose serverless solutions (BigQuery ML) when appropriate
    \item Build resilience into pipelines with retry and checkpoint mechanisms
    \item Monitor pipeline health and set up alerting
\end{enumerate}

\section{Further Reading}

\begin{itemize}
    \item Lakshmanan, V., et al. (2020). \textit{Machine Learning Design Patterns}. O'Reilly.
    \item Huyen, C. (2022). \textit{Designing Machine Learning Systems}. O'Reilly.
    \item Apache Airflow Documentation: \url{https://airflow.apache.org/docs/}
    \item Kubeflow Pipelines: \url{https://www.kubeflow.org/docs/components/pipelines/}
    \item MLflow Documentation: \url{https://mlflow.org/docs/latest/index.html}
    \item Optuna: \url{https://optuna.org/}
    \item BigQuery ML: \url{https://cloud.google.com/bigquery-ml/docs}
\end{itemize}

\section{Exercises}

\begin{exercise}
\textbf{End-to-End Training Pipeline}

Build a complete training pipeline using Airflow:

\begin{enumerate}
    \item Create a DAG with these tasks:
    \begin{itemize}
        \item Data extraction from BigQuery
        \item Data validation with Great Expectations
        \item Feature engineering
        \item Model training with MLflow tracking
        \item Model evaluation on holdout set
        \item Conditional model registration
    \end{itemize}
    \item Add error handling and retries
    \item Implement data quality gates
    \item Set up email alerting
    \item Create monitoring dashboard
\end{enumerate}
\end{exercise}

\begin{exercise}
\textbf{Hyperparameter Optimization}

Implement distributed HPO:

\begin{enumerate}
    \item Set up Optuna study with pruning
    \item Define search space for your model
    \item Implement objective function with MLflow integration
    \item Run 100+ trials with parallel execution
    \item Analyze parameter importance
    \item Visualize optimization history
    \item Train final model with best parameters
    \item Compare performance vs baseline
\end{enumerate}
\end{exercise}

\begin{exercise}
\textbf{Model Validation Framework}

Create comprehensive validation:

\begin{enumerate}
    \item Implement stratified k-fold cross-validation
    \item Calculate multiple metrics (accuracy, precision, recall, AUC)
    \item Analyze threshold-dependent metrics
    \item Generate confusion matrix and ROC curves
    \item Evaluate fairness across demographic groups
    \item Check for overfitting
    \item Create automated validation report
    \item Set performance thresholds for deployment
\end{enumerate}
\end{exercise}

\begin{exercise}
\textbf{A/B Test Design}

Design and analyze model A/B test:

\begin{enumerate}
    \item Define test hypothesis and metrics
    \item Calculate required sample size for 80\% power
    \item Implement hash-based user assignment
    \item Log predictions and outcomes
    \item Run statistical significance test
    \item Calculate confidence intervals
    \item Make deployment decision
    \item Document results and recommendation
\end{enumerate}
\end{exercise}

\begin{exercise}
\textbf{BigQuery ML Time Series}

Build production time series pipeline:

\begin{enumerate}
    \item Prepare time series data in BigQuery
    \item Train ARIMA\_PLUS model with seasonality
    \item Evaluate model performance
    \item Generate multi-step ahead forecasts
    \item Extract confidence intervals
    \item Implement anomaly detection
    \item Create Airflow DAG for daily retraining
    \item Validate forecast reasonableness
\end{enumerate}
\end{exercise}

\begin{exercise}
\textbf{Kubeflow Pipeline}

Implement Kubernetes-native ML pipeline:

\begin{enumerate}
    \item Define reusable components as Python functions
    \item Create data preprocessing component
    \item Build training component with GPU support
    \item Add evaluation component
    \item Implement deployment component
    \item Compile pipeline to YAML
    \item Run on Kubeflow cluster
    \item Monitor execution and debug failures
\end{enumerate}
\end{exercise}

\begin{exercise}
\textbf{Resilient Pipeline Design}

Build fault-tolerant pipeline:

\begin{enumerate}
    \item Implement retry with exponential backoff
    \item Add circuit breaker for external dependencies
    \item Create checkpoint mechanism
    \item Handle OOM errors with batch processing
    \item Implement graceful degradation
    \item Add dead letter queue for failures
    \item Set up comprehensive logging
    \item Test failure scenarios
\end{enumerate}
\end{exercise}
